% 宇宙学（综述）
% license CCBYSA3
% type Wiki

本文根据 CC-BY-SA 协议转载翻译自维基百科\href{https://en.wikipedia.org/wiki/Cosmology}{相关文章}。

宇宙学（Cosmology，源自古希腊语 κόσμος（cosmos，意为“宇宙”或“世界”）与 λογία（logia，意为“研究”））是一门研究宇宙——即宇宙整体本质的物理学与形而上学的分支。英语单词 “cosmology” 首次出现在 1656 年托马斯·布朗特（Thomas Blount）的《Glossographia》中，其含义为“关于世界的论述”。\(^\text{[2]}\)1731 年，德国哲学家克里斯蒂安·沃尔夫（Christian Wolff）在拉丁语中使用 “cosmologia” 一词，用以指代研究物质世界一般性质的形而上学分支。\(^\text{[3]}\)宗教或神话宇宙学（religious or mythological cosmology）是一套基于神话、宗教与秘传文献及传统的信仰体系，通常与创世神话（creation myths）及末世论（eschatology）相关。在天文学的科学范畴中，宇宙学主要关注对宇宙年代学（chronology of the universe）的研究。

物理宇宙学（physical cosmology）研究可观测宇宙的起源、其大尺度结构与动力学特征，以及宇宙的最终命运，同时探讨支配这些领域的科学定律。\(^\text{[4]}\)这一学科由科学家（如天文学家与物理学家）以及哲学家（如形而上学家、物理哲学家与时空哲学家）共同研究。由于其研究范围与哲学存在交叉，物理宇宙学中的理论既可能包含科学命题，也可能包含非科学命题，且可能依赖于无法通过实验检验的假设。物理宇宙学是天文学的一个子分支，专注于研究整个宇宙。现代物理宇宙学由“大爆炸理论”（Big Bang Theory）主导，该理论试图结合观测天文学与粒子物理学。\(^\text{[5][6]}\)更具体而言，标准的宇宙参数化模型包括暗物质（dark matter）与暗能量（dark energy），统称为“ΛCDM 模型”（Lambda-CDM model）。

理论天体物理学家戴维·N·斯珀格尔（David N. Spergel）将宇宙学称为一门“历史科学”（historical science），因为“当我们仰望宇宙时，实际上是在回望过去”，这源于光速的有限性。\(^\text{[7]}\)
\subsection{学科领域}
\begin{figure}[ht]
\centering
\includegraphics[width=8cm]{./figures/3be309e51e6413de.png}
\caption{} \label{fig_CoGy_1}
\end{figure}
物理学与天体物理学通过科学的观测与实验，在塑造人类对宇宙的理解过程中发挥了核心作用。物理宇宙学（physical cosmology）通过数学与观测相结合的方式，在对整个宇宙的分析中逐步形成。普遍认为，宇宙起源于“大爆炸”（Big Bang），随后几乎瞬间经历了“宇宙暴涨”（cosmic inflation）——即空间的迅速膨胀。据推测，宇宙大约在 $13.799 \pm 0.021$ 十亿年前由此诞生。\(^\text{[8]}\)宇宙起源学（cosmogony）研究宇宙的起源，而宇宙描绘学（cosmography）则负责绘制宇宙的结构与特征。

在狄德罗（Diderot）的《百科全书》（Encyclopédie）中，宇宙学被划分为四个部分：天学（uranology，研究天体之科学）、气学（aerology，研究空气之科学）、地学（geology，研究大陆之科学）以及水学（hydrology，研究水体之科学）。\(^\text{[9]}\)

形而上宇宙学（metaphysical cosmology）亦被描述为探讨人类在宇宙中与其他存在物之间关系的学科。罗马皇帝兼哲学家马可·奥勒留（Marcus Aurelius）对此作出了典型的论述，他指出人的位置与认知密切相关：“不知世界为何物者，不知己所处之地；不知世界为何而存者，不知己为何人，亦不知世界为何物。”\(^\text{[10]}\)
\subsection{发现}
\subsubsection{物理宇宙学}
物理宇宙学是物理学与天体物理学的一个分支，研究宇宙的物理起源与演化，同时也探讨宇宙在大尺度上的性质。在其最早期的形式中，它即如今所谓的“天体力学”（celestial mechanics）——对天体运动与“天界”的研究。希腊哲学家萨摩斯的阿里斯塔克（Aristarchus of Samos）、亚里士多德（Aristotle）与托勒密（Ptolemy）提出了不同的宇宙理论。其中，托勒密的地心体系（geocentric Ptolemaic system）长期占据主导地位，直到十六世纪，尼古拉·哥白尼（Nicolaus Copernicus）以及随后约翰内斯·开普勒（Johannes Kepler）与伽利略·伽利莱（Galileo Galilei）提出日心体系（heliocentric system）。这是物理宇宙学中最著名的认识论断裂（epistemological rupture）之一。

艾萨克·牛顿（Isaac Newton）于 1687 年发表的《自然哲学的数学原理》（Philosophiæ Naturalis Principia Mathematica）首次系统描述了万有引力定律（law of universal gravitation）。该理论为开普勒行星运动定律提供了物理机制，并能够解释行星之间引力相互作用所导致的系统异常。牛顿宇宙学与其前辈理论的一个根本区别在于哥白尼原理（Copernican principle）：地球上的物体遵循与天体相同的物理定律。这一思想构成了物理宇宙学的一个重要哲学飞跃。

现代科学宇宙学通常被认为始于 1917 年，当时阿尔伯特·爱因斯坦（Albert Einstein）发表了其对广义相对论（general relativity）的最终修正论文《广义相对论的宇宙学考量》（Cosmological Considerations of the General Theory of Relativity）。\(^\text{[11]}\)（尽管这篇论文直到第一次世界大战结束后才在德国以外被广泛传播。）广义相对论的提出促使威廉·德西特（Willem de Sitter）、卡尔·史瓦西（Karl Schwarzschild）与阿瑟·爱丁顿（Arthur Eddington）等宇宙起源学家（cosmogonists）探讨其在天文学中的影响，从而极大地提升了天文学家研究遥远天体的能力。物理学家们也由此开始放弃“宇宙是静止且不变的”这一传统假设。1922 年，亚历山大·弗里德曼（Alexander Friedmann）首次提出“宇宙膨胀”概念，即宇宙中包含运动的物质并在不断扩张。

与这种动态宇宙学观点并行的，是关于宇宙结构的一场长期争论——即 1917 年至 1922 年间的“大辩论”（Great Debate）。在这一时期，早期的宇宙学家如赫伯·柯蒂斯（Heber Curtis）与恩斯特·厄皮克（Ernst Öpik）通过观测确定，望远镜中所见的一些星云（nebulae）其实是距离我们极其遥远的独立星系。\(^\text{[12]}\) 柯蒂斯主张旋涡星云（spiral nebulae）本身就是独立的恒星系统，即“岛宇宙”（island universes）；而威尔逊山天文台的天文学家哈洛·沙普利（Harlow Shapley）则坚持认为宇宙仅由银河系这一恒星系统构成。两种观点的分歧在 1920 年 4 月 26 日于华盛顿特区美国国家科学院（U.S. National Academy of Sciences）会议上举行的“大辩论”中达到高潮。此争论最终在 1923 至 1924 年间得到解决，当时埃德温·哈勃（Edwin Hubble）在仙女座星系（Andromeda Galaxy）中探测到造父变星（Cepheid Variables）。\(^\text{[13][14]}\)其测得的距离证实了旋涡星云远远位于银河系边界之外。

此后，对宇宙的理论建模进一步探索了这样一种可能性：爱因斯坦在其 1917 年论文中引入的宇宙常数（cosmological constant），其数值可能导致宇宙处于膨胀状态。由此，比利时神父乔治·勒梅特（Georges Lemaître）于 1927 年提出了“大爆炸模型”（Big Bang model）。\(^\text{[15]}\)这一理论随后得到了哈勃 1929 年发现红移现象（redshift）的支持，\(^\text{[16]}\)并在 1964 年由阿诺·彭齐亚斯（Arno Penzias）与罗伯特·威尔逊（Robert Woodrow Wilson）发现宇宙微波背景辐射（cosmic microwave background radiation）后得到进一步验证。\(^\text{[17]}\)这些发现首次排除了多种其他竞争性的宇宙学模型。

自 1990 年前后起，观测宇宙学（observational cosmology）取得了多项重大突破，使宇宙学从主要依赖推测的学科转变为能够进行高精度预测的科学领域——理论与观测之间达成了前所未有的精确一致。这些进展包括：来自 COBE、\(^\text{[18]}\)WMAP、\(^\text{[19]}\)与 Planck\(^\text{[20]}\)卫星的宇宙微波背景观测；大型星系红移巡天（galaxy redshift surveys），如 2dFGRS\(^\text{[21]}\)与 SDSS\(^\text{[22]}\)；以及对遥远超新星与引力透镜（gravitational lensing）的观测。这些观测结果与“宇宙暴胀理论”（cosmic inflation theory）——即修正的大爆炸理论及其特定版本“ΛCDM 模型”——的预测高度吻合。因此，许多学者称当代为“宇宙学的黄金时代”（golden age of cosmology）。\(^\text{[23]}\)

2014 年，BICEP2 合作组宣称他们在宇宙微波背景中探测到了引力波（gravitational waves）的印记。然而，后续研究表明该结果并不可靠：所谓的引力波证据实际上是由星际尘埃造成的。\(^\text{[24][25]}\)

同年 12 月 1 日，在意大利费拉拉（Ferrara）召开的 Planck 2014 会议上，天文学家宣布：宇宙的年龄约为 138 亿年，其组成包括 4.9\% 的原子物质、26.6\% 的暗物质（dark matter），以及 68.5\% 的暗能量（dark energy）。\(^\text{[26]}\)
\subsubsection{宗教或神话宇宙学}
宗教或神话宇宙学是一套基于神话、宗教与秘传文献以及创世与末世传统（creation and eschatology）的信仰体系。几乎所有宗教都包含创世神话（creation myths），并且通常可依据米尔恰·伊利亚德（Mircea Eliade）及其同事查尔斯·朗（Charles Long）所建立的分类体系分为五种类型。

基于相似主题的创世神话类型（Types of Creation Myths based on similar motifs）包括：
\begin{enumerate}
    \item \textbf{无中生有的创造（Creation ex nihilo）}：创造源自神祇的思想、言语、梦境或身体分泌物。
    \item \textbf{潜水取地型创造（Earth diver creation）}：创造者派出潜水者（通常是鸟类或两栖动物）潜入原始之海的海底，从中取出沙土或泥浆，形成陆地世界。
    \item \textbf{逐层显现型创造（Emergence myths）}：始祖（progenitors）经历多个世界与形态的转变，直至抵达现今的世界。
    \item \textbf{肢解型创造（Creation by dismemberment）}：世界由原初存在（primordial being）的肢解与分裂而形成。
    \item \textbf{原初统一体分裂型创造（Creation by splitting or ordering of a primordial unity）}：宇宙由原始统一体（如“宇宙卵”的破裂或从混沌中建立秩序）而生。\(^\text{[27]}\)
\end{enumerate}
\subsubsection{哲学}
\begin{figure}[ht]
\centering
\includegraphics[width=6cm]{./figures/9ad29988d4f7e015.png}
\caption{以对数尺度表示的可观测宇宙示意图（Representation of the observable universe on a logarithmic scale）。从中心向外，距离由太阳开始逐渐增加。行星及其他天体的尺寸被放大，以便更清晰地展示其形状。} \label{fig_CoGy_2}
\end{figure}
宇宙学（Cosmology）研究作为整体的世界——包括空间、时间及一切现象。从历史上看，其研究范围极为广泛，并在许多情形下与宗教密切相关。\(^\text{[28]}\)某些关于宇宙的问题超出了科学探究的范围，但仍可通过其他哲学方法（如辩证法 dialectics）加以思考。超出科学范畴的探讨通常包括如下问题：\(^\text{[29][30]}\)
\begin{itemize}
    \item 宇宙的起源是什么？其“第一因”是否存在？其存在是否必然？（参见：一元论 monism、泛神论 pantheism、流溢论 emanationism、创世论 creationism）
    \item 宇宙的最终物质成分是什么？（参见：机械论 mechanism、动力论 dynamism、质形论 hylomorphism、原子论 atomism）
    \item 宇宙存在的最终理由（若有）是什么？宇宙是否具有目的？（参见：目的论 teleology）
    \item 意识的存在在现实的存在中是否起作用？我们如何得知自己关于宇宙整体的知识？宇宙学推理是否揭示形而上学真理？（参见：认识论 epistemology）
\end{itemize}
哲学史学者查尔斯·卡恩（Charles Kahn）将古希腊宇宙学的起源归因于阿那克西曼德（Anaximander）。\(^\text{[31]}\)





\subsection{历史宇宙学}
\begin{table}[ht]
\centering
\caption{}\label{tab_CoGy1}
\begin{tabular}{|c|c|c|c|}
\hline
\textbf{名称} & \textbf{作者与年代} & \textbf{分类} & \textbf{备注} \\
\hline
印度宇宙学 & 《梨俱吠陀》（Rigveda，约公元前 1700–1100 年）& 循环或振荡型（Cyclical or oscillating），时间上无限（Infinite in time） & 原初物质显现 311.04 万亿年后进入同等时长的非显现状态。宇宙显现 43.2 亿年后亦进入同等时长的非显现状态。无数个宇宙同时存在，这些循环自始至终永恒持续，由“欲望”（desires）所驱动。 \\
\hline
祆教宇宙学（Zoroastrian Cosmology）&《阿维斯陀经》（Avesta，约公元前 1500–600 年）& 二元宇宙论（Dualistic Cosmology） & 根据祆教宇宙学，宇宙是“存在与非存在”、“善与恶”、“光与暗”之间永恒冲突的显现。宇宙将在这种对立状态中持续 12000 年；在复活之时，这两种元素将再次被分离。\\
\hline
耆那教宇宙学（Jain cosmology）&《耆那教经典》（Jain Agamas，约公元 500 年撰写，依据大雄 Mahavira〔公元前 599–527 年〕的教义） & 循环或振荡型（Cyclical or oscillating），永恒而有限（eternal and finite） & 耆那教宇宙学认为“娑罗迦”（loka，即宇宙）是一个未被创造的永恒存在，其形状类似一个双腿分开、手臂支于腰间的站立之人。根据耆那教的描述，宇宙的上部宽阔，中部狭窄，而底部再次变得宽广。\\
\hline
巴比伦宇宙学（Babylonian cosmology）& 巴比伦文献（Babylonian literature，约公元前 2300–500 年）& 浮于无限“混沌之水”上的平面地球（Flat Earth floating in infinite "waters of chaos"）& 地球与天穹共同构成一个位于无限“混沌之水”（waters of chaos）中的整体；地球为平面且呈圆形，其上方覆盖着坚固的穹顶（即“天穹” firmament），以阻隔外部的“混沌之海”。\\
\hline
埃利亚学派宇宙学（Eleatic cosmology）& 巴门尼德（Parmenides，约公元前 515 年）  & 有限且球形（Finite and spherical in extent）& 宇宙是不变的、均匀的、完美的、必然的、超越时间的，既非生成也非消逝；虚空（Void）不存在。多样性与变化皆源于感官经验所导致的认知无知（epistemic ignorance）。时间与空间的界限皆为任意设定，且相对于巴门尼德式的整体（Parmenidean whole）而言是相对的。\\
\hline
数论派宇宙演化论（Samkhya Cosmic Evolution）& 迦毗罗（Kapila，公元前 6 世纪），阿修里（Asuri）之弟子 & “自性”（Prakriti，物质）与“神我”（Purusha，意识）之关系（Relation between Matter and Consciousness）& “自性”（Prakriti，物质）是生成世界（world of becoming）的根源，是一种纯粹的潜能（pure potentiality），依次演化为二十四种“原理”（tattvas）。这种演化之所以可能，是因为 Prakriti 始终处于其三种组成要素（称为“性” gunas）之间的张力状态之中：善性（Sattva，轻明或纯净）、动性（Rajas，激情或活动）、惰性（Tamas，惯性或沉重）。数论派的因果理论被称为“有因论”（Satkārya-vāda，theory of existent causes），其核心主张为：无物可凭空生灭，所有的演化不过是原始自然（primal Nature）从一种形式向另一种形式的转化。\\
\hline
圣经宇宙学（Biblical cosmology）&《创世纪》创世叙事（Genesis creation narrative）& 漂浮于无限“混沌之水”中的地球（Earth floating in infinite "waters of chaos"）& 地球与天穹构成一个位于无限“混沌之水”中的整体；其上方的“穹苍”（firmament）阻隔外部的“混沌之海”。\\
\hline
阿那克西曼德模型（Anaximander's model）& 阿那克西曼德（Anaximander，约公元前 560 年）& 地心说（Geocentric），圆柱形地球，空间无限而时间有限；首个纯机械模型（first purely mechanical model）& 地球静止地漂浮于无限空间的中心，不受任何支撑。\(^\text{[32]}\)起初，在冷热分离之后，一团火球出现，像树皮包裹树干一样环绕着地球。此火球破裂后形成了宇宙的其他部分。宇宙如同一个由火充盈的空心同心圆环系统，其边缘布满洞孔，如笛孔一般；并无真正的天体，唯有火光透过孔隙射出。共有三层圆环，从地球向外依次为：恒星（包括行星）、月亮与巨大的太阳。\(^\text{[33]}\)\\
\hline
原子论宇宙（Atomist universe）& 阿那克萨哥拉（Anaxagoras，公元前 500–428 年）及后来的伊壁鸠鲁（Epicurus）& 无限宇宙（Infinite in extent）& 宇宙中仅包含两种存在：无限数量的微小“种子”（即原子 atoms）与无限广阔的虚空（void）。所有原子由同一物质构成，但在大小与形状上有所差异。万物由原子的聚合而生，最终又分解为原子。此理论吸收了留基伯（Leucippus）的因果原理：“无事偶然，一切因理与必然而生。” 宇宙不受诸神统治。\(^\text{[citation needed]}\)\\
\hline
毕达哥拉斯宇宙（Pythagorean universe）& 腓罗劳斯（Philolaus，卒于公元前 390 年）& 存在“中央之火”（Central Fire）于宇宙中心（Existence of a “Central Fire” at the center of the Universe）& 在宇宙的中心存在一团“中央之火”，地球、太阳、月亮与行星皆以匀速围绕其旋转。太阳每年绕中央之火一周，而恒星则静止不动。地球在运行中始终以同一面朝向中央之火，因此该火永不为人所见。这是已知最早的“非地心宇宙模型”。\(^\text{[34]}\)\\
\hline
《论世界》（De Mundo）& 伪亚里士多德（Pseudo-Aristotle，卒于公元前 250 年，或约公元前 350–200 年间）& 宇宙是由天与地及其中的元素构成的系统（The Universe is a system made up of heaven and Earth and the elements contained within them）& 宇宙由“五种元素”组成，分布于五个球层中：地（earth）被水（water）包围，水被气（air）包围，气被火（fire）包围，而火又被以太（ether）所包围，这五层共同构成整个宇宙。\(^\text{[35]}\)\\
\hline
斯多亚派宇宙（Stoic universe）& 斯多亚学派（Stoics，公元前 300 年至公元 200 年）& 岛宇宙（Island universe）& 宇宙是有限的，被无限虚空（infinite void）所环绕。宇宙处于不断流变之中，其规模脉动变化，周期性地经历剧烈的动荡与宇宙火劫（conflagrations）。\\
\hline
柏拉图宇宙（Platonic universe） \quad 柏拉图（Plato，约公元前 360 年） \quad 地心说（Geocentric），复杂宇宙生成论（complex cosmogony），空间有限（finite extent），暗含有限时间（implied finite time），循环性（cyclical） \quad 静止的地球位于宇宙中心，其外依次环绕天体，这些天体在完美的圆周上运动，由造物主（\emph{Demiurge}）依意志所安排：月亮、太阳、行星与恒星。\(^\text{[35][][]}\)复杂的天体运动在每一个“完美之年”中循环重复。\(^\text{[39]}\)\\
\hline
欧多克索斯模型（Eudoxus' model） \quad 克尼多斯的欧多克索斯（Eudoxus of Cnidus，约公元前 340 年）及后来的卡利普斯（Callippus） \quad 地心说（Geocentric），首个几何—数学模型（first geometric-mathematical model） \quad 天体的运动如同附着在多个以地球为中心的同心不可见球体上，每个球体都围绕自身不同的轴线、以不同的速度旋转。\cite{ref40} 模型共设有二十七个同心球（homocentric spheres），每个球体用以解释特定天体的一类可观测运动。欧多克索斯强调这些球体纯为数学构造——它们在物理上并不存在，仅用于表征天体可能出现的位置。\cite{ref41}\\
\hline
亚里士多德宇宙（Aristotelian universe） \quad 亚里士多德（Aristotle，公元前 384–322 年） \quad 地心说（基于欧多克索斯模型），静态（static）、稳态（steady state）、空间有限（finite extent）、时间无限（infinite time） \quad 静止且呈球形的地球被 43 至 55 层同心天球所包围，这些天球为物质性的水晶结构。\cite{ref42} 宇宙自永恒以来保持不变。其体系包含第五元素——“以太”（\emph{aether}），作为对四种经典元素的补充。\cite{ref43}\\
\hline
阿里斯塔克斯宇宙（Aristarchean universe） \quad 阿里斯塔克斯（Aristarchus，约公元前 280 年） \quad 日心说（Heliocentric） \quad 地球每日自转一周，并以圆形轨道每年绕太阳公转一次。恒星球（sphere of fixed stars）以太阳为中心。\cite{ref44}\\
\hline
托勒密模型（Ptolemaic model） \quad 托勒密（Ptolemy，公元 2 世纪） \quad 地心说（Geocentric，基于亚里士多德宇宙） \quad 宇宙围绕静止的地球运转。行星在称为“本轮”（epicycle）的圆轨道上运动，而每个本轮的圆心又沿着更大的圆轨道（称为偏轮 eccentric 或均轮 deferent）绕地球附近的中心点运行。引入“均差点”（equants）的概念进一步提高了模型的复杂度，使天文学家得以精确预测行星位置。此模型因其长期稳定的预测能力，被认为是史上最成功的宇宙模型之一，其主要著作为《天文学大成》（\emph{Almagest}，又称“大体系” The Great System）。\\
\hline
* & * & * & * \\
\hline
* & * & * & * \\
\hline
* & * & * & * \\
\hline
* & * & * & * \\
\hline
* & * & * & * \\
\hline
* & * & * & * \\
\hline
* & * & * & * \\
\hline
* & * & * & * \\
\hline
* & * & * & * \\
\hline
* & * & * & * \\
\hline
* & * & * & * \\
\hline
* & * & * & * \\
\hline
* & * & * & * \\
\hline
* & * & * & * \\
\hline
* & * & * & * \\
\hline
* & * & * & * \\
\hline
* & * & * & * \\
\hline
* & * & * & * \\
\hline
* & * & * & * \\
\hline
* & * & * & * \\
\hline
* & * & * & * \\
\hline
* & * & * & * \\
\hline
* & * & * & * \\
\hline
* & * & * & * \\
\hline
* & * & * & * \\
\hline
* & * & * & * \\
\hline
* & * & * & * \\
\hline
* & * & * & * \\
\hline
* & * & * & * \\
\hline
* & * & * & * \\
\hline
* & * & * & * \\
\hline
* & * & * & * \\
\hline
* & * & * & * \\
\hline
* & * & * & * \\
\hline
* & * & * & * \\
\hline
* & * & * & * \\
\hline
\end{tabular}
\end{table}

